\section{Limitations of a software-only environment}\label{sec:limits}
The testbed exercises real 3GPP control-plane procedures and a real GTP-U user plane, so the
\emph{functional} conclusions transfer to a hardware deployment: correct slice selection,
subscription enforcement, per-slice UPF anchoring and the effect of the slice controls. The
\emph{quantitative} conclusions do not transfer directly, for the following reasons.

\begin{enumerate}
\item \textbf{No radio.} RLS replaces the entire NR stack below RRC/NAS. There is no PHY, no
  HARQ, no MAC scheduler, no numerology and no fading or interference. The air-interface
  latency (normally 0.5--4\,ms in the RAN user plane, depending on numerology and
  configured/dynamic grants) and its variability are absent. Slicing \emph{inside} the gNB
  (PRB partitioning, slice-aware scheduling) cannot be studied at all; our slice isolation
  is implemented in the core/transport only. Our sub-millisecond URLLC RTTs are therefore a
  lower bound on the processing cost of the core and transport, not a prediction of
  end-to-end URLLC latency.
\item \textbf{Shared, virtualised CPU.} Every NF, the gNB, all UEs, the traffic generators and
  the DN run on 4 vCPUs of one laptop SoC. Two artefacts follow and are quantified in this
  report. (i) CPU power management: an idle VM gives \emph{higher} RTT than a busy one
  (E7: p50 1.08\,ms idle vs.\ 0.10\,ms with pure busy-loops), so ``latency under load''
  comparisons must use the same CPU state. (ii) Scheduler noise: isolated per-second p99
  spikes of 5--27\,ms appear in every condition, including idle, and are caused by host or
  hypervisor scheduling, not by the 5G stack. We therefore compare medians and distributions
  over repeated runs rather than single worst cases.
\item \textbf{Throughput ceilings are CPU ceilings.} The 300--400\,Mbit/s eMBB maximum reflects
  user-space GTP-U processing in \code{nr-gnb}/\code{nr-ue} and \code{open5gs-upfd}, not any
  radio capacity. Absolute throughput numbers should not be extrapolated.
\item \textbf{Slice QoS enforcement.} Open5GS's UPF polices Session-AMBR (with up to about
  +15\,\% overshoot, \code{results/ambr\_check.txt}), but it does not schedule by 5QI or ARP.
  Differentiation by 5QI is therefore only signalled, not enforced. Our tc layer adds
  per-slice MBR and transport delay/loss, but it is a model of transport impairment, not of
  RAN scheduling.
\item \textbf{Clocks and measurement.} Only RTT is measured, because one-way delay would need
  synchronised clocks across processes. Per-second aggregation hides sub-second dynamics, and
  the 1-s KPI rows mean that the sandbox reacts with up to 2\,s of display lag.
\item \textbf{Scale.} Six UEs and 3000 virtual sensors behind two gateway UEs do not reproduce
  the RACH and signalling load of thousands of real devices. The mIoT result validates only
  the slice's user plane.
\item \textbf{One measured host.} Every number comes from the macOS host. A Windows PC runs
  the same software on a different CPU, hypervisor (Hyper-V) and kernel (Microsoft's WSL
  kernel), so its absolute latencies, throughput ceilings and CPU artefacts will differ. The
  Windows path is supported by the scripts. Its kernel prerequisites are checked by
  \code{wsl\_setup.sh}, which was itself validated in the Linux VM, but the experiments have
  not yet been repeated on a Windows PC.
\end{enumerate}

\noindent\textbf{Validity of conclusions.} The claims that \emph{do} hold:
(a) the end-to-end 5G procedures and slice selection work as specified;
(b) a per-slice control (MBR, delay) is applied within one KPI interval and affects only its
slice;
(c) with eMBB held at its MBR, URLLC and mIoT are unaffected by an eMBB surge;
(d) under platform overload, moving URLLC onto the eMBB UPF shifts its latency distribution
upwards. The per-run effect is consistent in direction but not yet statistically
significant with three runs per mode (\S\ref{sec:e5}).
Claims about absolute latency budgets against TS~22.104 targets would need a hardware or
PHY-emulating RAN (e.g.\ OAI/srsRAN with an SDR, or OAI's RF simulator) and a dedicated core host.
